\documentclass[12pt,a4paper]{article}
\usepackage{amsmath,amssymb,amsthm,geometry,hyperref}
\geometry{margin=1in}
\title{Laakkonen's Paradox: A Constructive Alternative to the Classical Metric of the Origin}
\author{Rob Laakkonen --- Infinite Space Mechanic \\ Language Framework: Lean 4 Theorem Prover \\ GitHub: Infinitespacemechanic/Laakkonen-Paradox}
\date{2026}
\begin{document}
\maketitle

\begin{abstract}
Classical field theories create mathematical singularities at geometric coordinates by conflating spatial indices with physical material states. By permitting volume and mass boundaries to drop to an absolute zero ($0$), classical physics creates a structural paradox at the core of massive bodies where continuous downward mechanical pressure (overburden) meets an unphysical null field. This paper presents Laakkonen's Paradox and its computational resolution: a constructive framework where the ground state of physical measures is structurally locked to a universal unit hold ($1$). We present the complete physical background, the decoupling of vector cancellation from field intensity, and a rigorous, computer-checked mathematical type system in Lean 4 that completely prevents structural collapse to zero.
\end{abstract}

\section{The Paradox of the Origin: Label vs. State}
The foundational error in textbook Newtonian and relativistic physics lies in a structural equivocation between a geometric coordinate label and a physical state value.

Classical mechanics describes the gravitational acceleration $g(r)$ inside a uniform massive sphere as:
$$ g(r) = \frac{G \cdot M(r)}{r^2} $$

As the coordinate radius approaches the absolute origin ($\lim_{r\to 0}$), classical models allow the enclosed volume and consequently the enclosed mass $M(r)$ to shrink smoothly to absolute zero.

The Mechanical Breakdown: This derivation treats geometric cancellation as a material void. At the center of a planetary body, a ``qazillion pounds'' of overburden is continuously pressing down. If gravity were truly zero as a state value at the center, this immense overburden would mechanically float. The coordinate $r=0$ is merely the address of the origin. The physical state occupied by that address cannot be nothing; it is an active, locked rest-state.

\section{Resolving the Field Metrics}
To resolve, separate Net Vector Acceleration $\vec{g}_{net}$ from Field Potential Intensity $g_{field}$.

\begin{itemize}
\item Balanced Vectors: $\vec{g}_{net}(0) = \vec{0}$
\item Maximum Field Intensity: $g_{field}(0) = g_{max} \ge 1$
\end{itemize}

Mass cannot vanish into nothingness at the core, nor can it dissipate into a true absolute vacuum. Space is continuous and continuously occupied.

$$\text{No Zero} \rightarrow \text{Greater Than Zero} \rightarrow \text{System Holds At 1 Until Exact}$$

\section{The Radial Field Continuum}
\begin{enumerate}
\item Origin Center ($r=0$): Locked at $M_{hold}=1$
\item Inner Iron Zone: High-density accumulation under active grav lock ($M(r) \gg 1$)
\item Middle Slot: Resonant containment regions (rings, asteroid belts, Van Allen)
\item Gaseous Envelope: Diffuse distributions where pressure accrues continuously
\item Dry Space: Outer boundary traditionally called vacuum, reclassified as dry space ($\lim_{r\to\infty} M(r)=1$)
\end{enumerate}

\section{Formal Type-Theoretic Foundations in Lean 4}

\subsection{Construction of the Non-Zero Field (LaakReal)}
\begin{verbatim}
@[ext]
structure LaakReal where
  val : R
  prop : val >= 1
\end{verbatim}

Zero is unconstructible:
\begin{verbatim}
theorem LaakReal.no_zero_exists : ¬ ∃ (z : LaakReal), z.val = 0
\end{verbatim}

\subsection{Translation Algebras}
Shifted Addition: $a \oplus b = a+b-1$, Multiplication: $a \otimes b = a\cdot b$, where $1$ is identity.

\subsection{Bounded Saturated Calculus}
If $M(r) \to 1$ then $dM/dr \to 0$. Field flattens asymptotically.

\begin{verbatim}
def LaakReal.sub (a b : LaakReal) : LaakReal where
  val := if a.val - b.val + 1 >= 1 then a.val - b.val + 1 else 1

theorem LaakReal.sub_self (a : LaakReal) : a - a = LaakReal.one
theorem LaakReal.derivative_step_never_zero (a b : LaakReal) : (a - b).val != 0
\end{verbatim}

\section{Conclusion}
Laakkonen's Paradox exposes the compromise of classical physics: treating coordinate label $r=0$ as material value $0$. By introducing a computer-verified, non-vanishing unit space (LaakReal), this framework provides a robust alternative for modeling continuously occupied physical fields. The universe is protected from structural null singularities at its dense cores and empty boundaries alike; the system holds at 1 until exact.

\end{document}
